\title{Byte Latent Transformer: Patches Scale Better Than Tokens}

\begin{abstract}
We present the Byte Latent Transformer ({\textsc BLT}), a novel large language model operating at the byte level that, for the first time, achieves parity with token-based LLMs in large-scale settings, while offering notable gains in both inference efficiency and model robustness. In {\textsc BLT}, bytes are aggregated into adaptively sized patches, which function as the core computational units. Patch boundaries are determined via the entropy of forthcoming bytes, thereby directing more computational effort and network capacity toward regions with greater data complexity. We conduct the first flop-controlled scaling analysis for byte-level architectures, exploring models with up to 8B parameters trained on 4T bytes. Our findings confirm that it is practical to scale raw byte-based models without needing a predefined vocabulary. Dynamic selection of extended patches over predictable data segments enhances efficiency in both training and inference, while yielding qualitative improvements in tasks involving reasoning and generalization to rare cases. In summary, at equivalent inference budgets, {\textsc BLT} achieves superior scaling compared to models relying on tokenization, simultaneously increasing both patch size and model capacity.
\end{abstract}

\section{Introduction}
\label{section:intro}

We present the Byte Latent Transformer~(\textbf{BLT}), an architecture that operates directly on raw byte sequences without relying on tokenization. For the first time, BLT matches the large-scale performance of tokenization-dependent models, while achieving notable gains in both efficiency and robustness (\S\ref{sec:robustness}). Conventional large language models (llms) typically involve an almost completely end-to-end training paradigm, with the exception of tokenization—a rule-based pre-processing approach that transforms byte sequences into a predefined, discrete set of tokens. This tokenization process introduces compression biases, manifesting as weaknesses such as increased sensitivity to specific domains or modalities~\citep{dagangetting}, greater vulnerability to input perturbations~(\S\ref{sec:robustness}), insufficient awareness of orthographic structure~\citep{edman2024cute}, and uneven support across languages~\citep{liang2023xlm,petrov2024language,limisiewicz2024myte}.

Historically, tokenization has played a pivotal role because training large language models on raw bytes is infeasible at scale, given the prohibitively lengthy input sequences~\citep{xue2022byt5}. Some studies have sought to address this by introducing more resource-efficient self-attention mechanisms~\citep{el2020characterbert, clark2022canine} or opting for architectures that forgo attention altogether~\citep{wang2024mambabyte}~(\S\ref{section:related_work}). Nonetheless, these solutions are most effective for \textit{small models}. In large-scale settings, the majority of computation within Transformers is consumed by extensive feed-forward network layers applied to every byte, rather than by the attention component.

To optimize the distribution of computational resources, we introduce an adaptive and trainable approach for aggregating bytes into \textit{patches}~(\S\ref{section:patching}), as well as a novel model architecture that integrates information from both bytes and patches. In contrast to tokenization, {\textsc BLT} does not utilize a predetermined patch vocabulary. Instead, any collection of bytes can be transformed into latent patch representations using compact, learned encoder and decoder modules. Our findings demonstrate that this strategy achieves a \textit{higher} level of compute efficiency compared to models built on tokenization.

Traditional tokenization-based large language models utilize an equal amount of computation for each token, prioritizing consistency at the cost of efficiency. The compression strategies underlying token induction do not always align well with the actual complexity of prediction tasks. Our approach centers on the principle that computational resources should be distributed adaptively, depending on local needs. For instance, predicting common word endings typically requires minimal computation, as these constitute straightforward, low-entropy tasks, whereas generating the initial word in a sentence demands more substantial modeling power. Reflecting this insight, {\textsc BLT}'s architecture (\S\ref{section:architecture}) comprises three transformer components: two compact, byte-level \textit{local models}, and a much larger, global \textit{latent transformer}~(\autoref{fig:arch_high_level}). 

To facilitate this adaptive computation, {\textsc BLT} organizes input bytes into patches according to the entropy associated with the next-byte prediction. This method yields context-sensitive groupings of bytes, each with a similar density of information, guiding the dynamic allocation of computation throughout the model.

We conduct the first comprehensive flop-controlled scaling analysis of byte-level models at scales reaching 8B parameters and 4T training bytes, demonstrating the feasibility of training models directly on raw byte sequences at scale, entirely omitting fixed-vocabulary tokenization. Our findings indicate that {\textsc BLT} achieves comparable performance to Llama 3 under equivalent training flop constraints\footnote{Model computational demands are measured in Floating Point OPerations (flops).}, while delivering up to 50\% inference flop savings (\S\ref{section:scaling}). 

Additionally, we find that leveraging raw bytes yields notable gains when modeling rare and atypical data. The {\textsc BLT} models show increased resilience compared to tokenization-based approaches in the face of noisy input, and they exhibit improved handling of character-level tasks, including orthographic analysis, phonological representation, and machine translation in low-resource settings~(\S\ref{sec:robustness}).

Moreover, with {\textsc BLT}, both model and patch sizes can be scaled up simultaneously without surpassing a fixed inference flop threshold. Enlarging patch sizes tends to reduce computational cost on average, freeing up resources to further expand the global latent transformer, as it is invoked less frequently. Our inference-flop controlled scaling evaluations (\autoref{fig:fixed_inference_scaling}) reveal much more favorable scaling behavior relative to token-based model architectures.

To conclude, this work offers several key contributions: 1) We present {\textsc BLT}, a novel byte latent LLM framework capable of adaptively distributing computational resources to enhance flop efficiency; 2) We demonstrate that our method reaches flop-controlled training parity with Llama 3 up to the 8B parameter range, while permitting substantial flop efficiency improvements—up to 50\%—with only minimal impacts on evaluation scores; 3) {\textsc BLT} models open a new pathway for scaling large language models, enabling the expansion of model size without exceeding a constant inference-time computational cost; 4) Our findings establish that {\textsc BLT} exhibits greater resilience to input perturbations, and that its design captures sub-word input features that conventional token-based LLMs overlook. Training and inference implementations for {\textsc BLT} are available at~\url{https://github.com/facebookresearch/blt}.

\section{Patching: From Individual Bytes to Groups of Bytes}
\label{section:patching}

Dividing bytes into \textit{patches} enables {\textsc BLT} to assign computational resources in an adaptive manner based on the data’s context. Figure~\ref{fig:patching_types} illustrates multiple approaches for splitting bytes into patches. More precisely, a patching function $f_p$ partitions a byte sequence $\pmb{x}=\{x_i,|i=1,\ldots n\}$ of length $n$ into $m < n$ patches $\pmb{p}=\{p_j|j=1,\ldots,m\}$, where each $x_i$ is mapped to the set \{0,1\}, and a value of 1 denotes the initiation of a new patch. For both models based on tokens and those based on patches, the principal computational expense of handling input is governed by the count of steps performed by the core Transformer module. In {\textsc BLT}, this translates to the number of patches that must be generated by the specific patching function used. As a result, the mean length of a patch—or the \textit{patch size}—serves as the primary determinant of the computational effort required to process data at both training and inference time with a chosen patching scheme~(\S\ref{section:flops}).

We next outline three distinct patching functions: fixed-length byte patching~(\S\ref{section:static-patch}), patching at whitespace~(\S\ref{section:white-patch}), and adaptive patching based on local entropies produced by a small byte-level language model~(\S\ref{section:dyn-patch}). We also elaborate on incremental patching and contrast patching with tokenization~(\S\ref{section:bpe-patch}).

\subsection{Strided Patching Every K Bytes}
\label{section:static-patch}

A common approach to organizing bytes involves partitioning them into uniform patches of size $k$, as implemented in MegaByte~\citep{yu2023megabyte}. This method, using a constant stride, offers implementation simplicity for both training and inference, and enables direct adjustment of the average patch size, thus making the computational requirements easily tunable. Nevertheless, there are notable limitations to this technique. Primarily, the allocation of computational effort is fixed rather than adaptive: this rigidity may result in inefficiencies, such as a transformer step $j$ being expended merely to process uninformative regions like whitespace in code, or in contrast, providing inadequate attention to more information-rich bytes, such as those found in mathematical content. Additionally, fixed-size patching can create non-uniform and non-context-sensitive groupings of similar byte sequences; for example, identical words might be divided differently depending on their position, leading to inconsistent segmentation.

\subsection{Space Patching}
\label{section:white-patch}

\citet{slagle2024spacebyte} introduces an intuitive yet powerful modification to strided patching by generating new patches immediately following any space-like byte\footnote{
    Space-like bytes refer to any bytes excluding Latin letters, digits, or utf-8 continuation bytes. Furthermore, every patch must include at least one non space-like byte.
}—positions that often correspond to natural linguistic boundaries across multiple languages. With Space patching, each word is given a dedicated latent transformer computation step (i.e., more flops), ensuring consistent patching of words throughout sequences and targeting additional computation toward challenging predictions, which typically arise right after spaces. For instance, when answering “Who composed the Magic Flute?\uline{\hspace{.6em}}”, determining the first byte of the answer involves far greater ambiguity than producing subsequent bytes after the initial “M,” since the initial character drastically restricts the pool of plausible responses—thereby making it much easier to generate the rest of the answer “Mozart.” Nonetheless, Space patching exhibits limitations: it does not generalize seamlessly to all languages and domains, and it is unable to flexibly adjust patch size. In the following section, we introduce a novel patching technique grounded in the observation that predicting the first bytes of words tends to be the hardest, while also offering a principled approach to dynamically controlling patch size.

\subsection{Entropy Patching: Using Next-Byte Entropies from a Small Byte LM}
\label{section:dyn-patch}

Instead of employing rule-based heuristics like whitespace for boundary detection, we adopt a data-centric strategy aimed at pinpointing next-byte predictions with high uncertainty. Specifically, we propose \textit{entropy patching}, a method that establishes patch boundaries by utilizing entropy estimates.

We train a small byte-level auto-regressive language model on the training data for {\textsc BLT} and compute next byte entropies under the LM distribution $p_{e}$ over the byte vocabulary $\mathcal{V}$:

\begin{equation}
H(x_i) = \sum_{v \in \mathcal{V}} p_{e}(x_i=v|\pmb{x}_{<i}) \log p_{e}(x_i=v|\pmb{x}_{<i})
\end{equation}

We explore two approaches for detecting patch boundaries based on entropies $H(x_i)$. The first method selects points that surpass a global entropy threshold, as shown in~\autoref{fig:patching}. The second method targets points whose entropy is significantly higher compared to the preceding value. This latter strategy can also be understood as locating instances where the roughly monotonic decrease in entropy within a patch is disrupted.

\begin{align*}
    \text{Global Constraint}\hspace{1cm}H(x_t)& > \theta_g\\
    \text{Approx. Monotonic Constraint}\hspace{1cm}H(x_t)& - H(x_{t-1}) > \theta_r
\end{align*}

Patch boundaries are determined through a lightweight preprocessing procedure applied at the dataloading stage. This approach contrasts with~\citet{nawrot-etal-2023-efficient}, where a classifier is trained to detect patch boundaries based on entropy estimates. In our experimental setup~(\S\ref{section:experiments}), we directly compare these two techniques for differentiating between bytes with low and high entropy.

\subsection{The Byte-Pair Encoding (BPE) Tokenizer and Incremental Patching}
\label{section:incremental-patching}
\label{section:bpe-patch}

Numerous contemporary llms, such as the baseline Llama 3, implement subword tokenization methods like bpe~\citep{gage1994new, sennrich-etal-2016-neural}. In this context, we define ``tokens'' as byte sequences selected from a \textit{finite} set specified before training, distinguishing them from ``patches,'' which are adaptively grouped segments that do not depend on a predetermined vocabulary. An essential distinction between tokens and patches lies in the fact that, for tokens, the underlying byte-level information is inaccessible to the model.

An important advancement introduced by {\textsc BLT} over models based on traditional tokenization is its reconfiguration of the balance between vocabulary size and computational resources. Conventional large language models (LLMs) must contend with the fact that enlarging the vocabulary typically results in bigger average token lengths, thus allowing the model to complete sequences in fewer steps. However, this also inflates the dimensionality of the output in the model’s final projection layer. This balance constrains tokenization-based strategies from meaningfully adjusting token lengths or reducing inference costs. To illustrate, Llama 3 achieves an increase in average token size—from 3.7 to 4.4 bytes—but at the expense of expanding its embedding table by a factor of four compared to Llama 2.

During generation, {\textsc BLT} must determine at each position in the byte sequence whether it lies at the boundary of a patch, since this decision affects whether the Latent Transformer is applied. Importantly, this determination must rely solely on information available up to the current step, independent of any future bytes which have not yet been produced. Therefore, the segmentation mechanism cannot depend on knowledge of subsequent parts of the sequence. Formally, a patching function $f_p$ is said to support incremental patching if the following condition is met:
$$f_p(\pmb{x}_{<i}) = f_p(\pmb{x})_{<i}$$
It should be noted that bpe does not meet the requirements for an incremental patching scheme, as the way a prefix is tokenized may vary according to the context provided by the remainder of the sequence. Consequently, it fails to fulfill the condition described above\footnote{Introducing a dedicated delimiter token to mark patch boundaries can enable bpe to operate as an incremental patching scheme, though this comes at the cost of lengthening the resulting byte sequence.}.

\section{{\textsc BLT} Architecture}
\label{section:architecture}
The {\textsc BLT} model architecture consists of a large global autoregressive language model that processes inputs at the patch level. In addition, two lightweight local models are employed: one converts sequences of bytes into patch representations, while the other reconstructs bytes from the patch-level outputs~(\autoref{fig:arch_high_level}).

\subsection{Latent Global Transformer Model}

\textit{The Latent Global Transformer} refers to an autoregressive transformer architecture, denoted as $\mathcal{G}$, with $l_{\mathcal{G}}$ layers, designed to translate a series of input latent patch representations $p_j$ into corresponding output patch representations $o_j$. In this manuscript, the notation $j$ is consistently employed for patch indices, whereas $i$ is reserved for indexing bytes. The global component applies a block-causal attention mask~\citep{dubey2024llama}, ensuring that each position may only attend to itself and preceding patches within a given document. This model is responsible for the majority of computational operations, both in pre-training and inference stages, making the selection of its application points crucial for modulating and distributing computational resources according to the complexity of different input and output segments.

\subsection{Local Encoder}
\label{section:local}

\textit{The Local Encoder Model}, represented by $\mathcal{E}$, employs a streamlined transformer architecture featuring $l_{\mathcal{E}} << l_{\mathcal{G}}$ layers. Its principal function is to efficiently convert input byte sequences $b_i$ into highly descriptive patch representations $p_j$. Unlike traditional transformer models, this design introduces a cross-attention mechanism subsequent to each transformer layer, allowing for the pooling of byte-level embeddings into patch-level representations~(\autoref{fig:crossattn}).

The initial phase embeds the input byte sequence $b_i$ into continuous vectors using an embedding matrix of dimension $\mathbb{R}^{256 \times h_{\mathcal{E}}}$, with these outputs referred to as $x_i$. Optionally, these embeddings are enriched with supplementary hash-embeddings as discussed in~(\S\ref{sec:hash-emb}). 

The input representations are further refined through a repeated sequence of transformer and cross-attention layers, as outlined in~(\S\ref{sec:enc-cross-attn}), eventually yielding patch representations $p_i$, which are subsequently handled by the global transformer $\mathcal{G}$. Within each transformer layer, a \textit{local block causal} attention mask is employed; here, each byte can access information from a prescribed window of $w_{\mathcal{E}}$ preceding bytes, potentially spanning across moving patch boundaries, but always constrained by document boundaries. The upcoming subsections elaborate on the embedding strategy and cross-attention mechanism in greater detail.

\subsubsection{Encoder Hash n-gram Embeddings}
\label{sec:ngram-emb}
\label{sec:hash-emb}
An essential aspect of constructing resilient and expressive representations at every position $i$ involves integrating data regarding the bytes that come before. Within {\textsc BLT}, this is accomplished by representing the current byte $b_i$ on its own, as well as embedding it within a broader context as a byte n-gram. At each index $i$, our process begins by forming byte-grams

\begin{align}
    g_{i,n}=\{b_{i-n + 1},\ldots, b_i\}
\end{align}

for each byte position $i$ and for values of $n$ ranging from three to eight.\footnote{
Byte-grams are excluded for sizes $n$ or greater when $i < n$.}

Subsequently, we propose hash-based $n$-gram embeddings, wherein every byte $n$-gram is transformed using a hash function into an index within a corresponding embedding table $E_{n}^{hash}$ of fixed dimensionality, for each $n$ in the set $\{3,4,5,6,7,8\}$~\citep{bai2010learning}.
The embeddings obtained from this mapping are then added to the original byte’s embedding, followed by normalization, before serving as input to the local encoder. The augmented embedding is computed as follows:

\begin{align}
e_i &= x_i + \sum_{n = 3,...,8}E_{n}^{hash}(\text{Hash}(g_{i,n})) \\
\text{where, Hash}(g_{i,n}) &= \text{RollPolyHash}(g_{i,n}) \% |E_{n}^{hash}|
\end{align}

We divide $e_i$ by the sum of the number of $n$-gram sizes and one for normalization, utilizing RollPolyHash as outlined in Appendix~\ref{appendix:rollpolyhash}. Section~\ref{section:ablations} investigates how varying both the value of $n$ and the embedding table size for $n$-gram hash embeddings influences scaling law trends under fixed FLOP conditions. Alongside experiments with hash $n$-gram embeddings, we also evaluated frequency-based $n$-gram embeddings; a comprehensive account of these trials is included in Appendix~\ref{appendix:ngrams}.

\subsubsection{Encoder Multi-Headed Cross-Attention}
\label{sec:enc-cross-attn}
Our approach builds upon the input cross-attention component introduced in the Perceiver framework~\citep{jaegle2021perceiver}. A key distinction in our method is that, rather than relying on a predetermined collection of latent representations, we associate each latent with a specific patch, as depicted in~(\autoref{fig:crossattn}). Each patch representation focuses solely on the bytes contained within its patch. For each patch $p_j$, a query vector is computed by first aggregating (e.g., pooling) its constituent byte-level representations, which is subsequently transformed through a linear layer, $\mathcal{E}_{C} \in \mathbb{R}^{h_{\mathcal{E}} \times (h_{\mathcal{E}} \times U_{\mathcal{E}})}$, where $U_{\mathcal{E}}$ defines the total number of cross-attention heads in the encoder. Specifically, let $f_{\text{bytes}}(p_j)$ represent the byte sequence that forms patch $p_j$; we then proceed to compute

\begin{align}
P_{0,j} &= \mathcal{E}_{C}(f_\text{bytes}((p_j)), f~ \text{is a pooling function} \\
P_l &= P_{l-1} + W_o\left(\text{softmax}\left(\frac{QK^T}{\sqrt{d_k}}\right)V\right) \\
\text{where } Q_j &= W_q(P_{l-1,j}), K_i = W_k(h_{l-1,i}), V_i = W_v(h_{l-1,i}) \\
h_l &= \text{Encoder-Transformer-Layer}_l(h_{l-1})
\end{align}

In this context, $P \in \mathbb{R}^{n_p \times h_{\mathcal{G}}}$ denotes the collection of $n_p$ patch representations, which are fed into the global model. The initial values for $P$ are obtained by aggregating the byte embeddings $e_i$ corresponding to the bytes within each patch $p_j$. The matrices $W_q$, $W_k$, $W_v$, and $W_o$ serve as the projection weights for generating queries, keys, values, and outputs, respectively; specifically, the keys and values are obtained by projecting byte-level representations $h_i$ from the prior layer (or from $e_i$ in the first layer). To ensure that attention within a patch is localized, we apply a tailored masking mechanism such that a given query $Q_j$ can only interact with keys and values associated with bytes in its own patch $j$. Since we employ a multi-head attention mechanism over $Q$, $K$, and $V$, and since patch representations commonly have a larger dimensionality $h_{\mathcal{G}}$ compared to $h_{\mathcal{E}}$, we organize $P_l$ as multiple attention heads each of size $h_{\mathcal{E}}$ during the cross-attention operation; these heads are then concatenated to form a vector of dimension $h_{\mathcal{G}}$. Furthermore, we employ pre-LayerNorm normalization to the queries, keys, and values, and refrain from using positional embeddings in this cross-attention layer. A residual connection is also introduced around the cross-attention block to facilitate optimization.

\subsection{Local Decoder} 

Like the local encoder, the local decoder $\mathcal{D}$ utilizes a compact transformer architecture, comprising $l_{\mathcal{D}} << l_{\mathcal{G}}$ layers. Its task is to reconstruct a series of global patch embeddings $o_j$ back into raw bytes $y_i$. During this process, the local decoder generates the byte sequence by conditioning each prediction on the bytes previously decoded. It does so by consuming the internal states output by the local encoder for the byte sequence. The architecture alternates between $l_{\mathcal{D}}$ layers of cross-attention and standard transformer blocks. Each cross-attention layer in the decoder precedes a transformer layer, transforming patch embeddings into byte-level representations; the subsequent transformer layer in the local decoder then operates over this constructed byte sequence.

\subsubsection{Decoder Multi-headed Cross-Attention} 

In the decoder cross-attention, the roles of the queries and key/values are interchanged i.e. the byte-representations are now the queries, and the patch representations are now the key/values. The initial byte-representations for the cross-attention are initialized as the byte embeddings from the last encoder layer i.e. $h_{l_{\mathcal{E}}}$. The subsequent byte-representations for layer $l$, $d_{l,i}$ are computed as:

\begin{align}
D_0 &= h_{l_{\mathcal{E}}} \\
B_l &= D_{l-1} + W_o\left(\text{softmax}\left(\frac{QK^T}{\sqrt{d_k}}\right)V\right), \\
  &\text{where } Q_i = W_q(d_{l-1,i}), K_i = W_k(\mathcal{D}_C(o_j)), V_i = W_v(\mathcal{D}_C(o_j)) \\
  D_l &= \text{Decoder-Transformer-layer}_l(B_l)
\end{align}

In this context, $W_k$ and $W_v$ refer to the key and value projection matrices, which act through a linear transformation and splitting operation $\mathcal{D}_C$ on the global model's final patch representations $o_j$. The query projection matrix $W_q$ is applied to the byte-level inputs $d_{l-1}$ from the previous decoder transformer layer (using $h_{l_{\mathcal{E}}}$ at the initial layer), while $W_o$ denotes the output projection matrix. Consequently, $B$ assumes a shape of $\mathbb{R}^{h_{\mathcal{D}} \times n_b}$, where $n_b$ designates the number of output bytes. The subsequent decoder states, $D_l$, result from processing the cross-attention block output, $B$, through a decoder transformer layer. Consistent with the local encoder's cross-attention mechanism, we employ multi-head attention, utilize pre LayerNorms, exclude positional embeddings, and implement a residual connection that wraps the cross-attention component.

\section{Experimental Setup}
\label{section:experiments}
We establish rigorously controlled experiments to evaluate {\textsc BLT} against tokenization-based approaches, ensuring that {\textsc BLT} is not granted any potential benefits arising from access to longer sequence lengths.

\subsection{Pre-training Datasets} The models evaluated in this study are pre-trained utilizing two primary datasets: (1) The Llama 2 dataset~\citep{touvron2023llama}, which contains 2 trillion tokens drawn from an array of publicly available sources. This corpus undergoes extensive cleaning and filtering procedures to enhance its overall quality; (2) {\textsc BLT}-1T, a newly constructed dataset consisting of 1 trillion tokens obtained from a mix of public datasets, augmented by a subset derived from the Datacomp-LM pre-training collection~\citep{li2024datacomp}. For scaling law analyses aimed at determining optimal token counts as described in~\cite{dubey2024llama}, the first dataset serves as the primary source for investigating ideal architectural design for {\textsc BLT}. The second dataset, {\textsc BLT}-1T, is dedicated to end-to-end pre-training trials, enabling direct comparisons with Llama 3 for various downstream evaluation tasks. Importantly, neither dataset incorporates information sourced from Meta platforms or services. Additionally, to standardize baseline evaluations for tokenizer-dependent models, experiments employ the Llama 3 tokenizer, configured with a 128K-token vocabulary. This choice is motivated by observed performance gains relative to models built using the Llama 2 tokenizer.

\subsection{Entropy Model} For our experiments, the entropy model is implemented as a byte-level language model, trained using the identical data distribution as the comprehensive {\textsc BLT} model. By default, the architecture consists of a transformer with 100M parameters, comprising 14 layers and employing a hidden size of 512. The model utilizes sliding window attention spanning 512 bytes. All other hyperparameter choices align with those used in our local and global transformer models. As detailed in \autoref{section:ablations}, we evaluated a range of model capacities, context window sizes, and architectural variants. Notably, if the receptive field is constrained to be sufficiently limited, the learned entropy model can be represented as a compact and efficient lookup table.

\subsection{Entropy Threshold and Equalizing Context Length} For entropy-based patching approaches, we determine a patching threshold designed to yield a specified mean \textit{patch size} when evaluated over the pretraining data mixture. In contrast to tokenization, {\textsc BLT} permits the selection of \textit{patch size} to be fully flexible, which can substantially affect the total context consumed by the model. To prevent models with larger patch sizes from benefiting from an extended context window, we keep the expected total number of bytes per batch constant, thereby normalizing context length across configurations. Consequently, models employing larger patch sizes use proportionally shorter sequences. Specifically, for Llama 2 data, an 8k byte context is adopted, while for the {\textsc BLT}-1T dataset, the context window is increased to 16k bytes on average, all while holding the mean batch size fixed at 16M bytes.

Although the average batch size remains unchanged, dynamic patching techniques lead to varying proportions of bytes per patch when assembling batches. In our {\textsc BLT} training implementation, we group patches efficiently within each batch to minimize the need for padding operations in the computationally intensive latent transformer. This design maintains a consistent patch count across all batches. To prevent memory overflow caused by exceptionally large patches, we pad and, if necessary, truncate input byte sequences to 12k and 24k bytes, respectively, for the Llama 2 and {\textsc BLT}-1T datasets during training.

\subsection{Entropy Model Context}
\label{section:entropy-context}
Our empirical analysis indicates that applying entropy patching tends to result in increasingly larger patches for structured tasks such as multiple choice problems (refer to \autoref{fig:mmlu_patching}), which often contain considerable repetition. These differences stem from the entropy model context exhibiting reduced entropy in these recurring segments. Therefore, for the large-scale implementation of {\textsc BLT}-Entropy with a patch size of 4.5, we refresh the entropy context by inserting new lines and implement an approximate monotonicity constraint; this method is more robust to "entropy drift" brought about by variations in context length. Notably, this adjustment only modifies the way entropies are computed, and we maintain the same approach for determining the appropriate entropy threshold value.

\subsection{FLOPs Estimation} 
\label{section:flops}

Our approach for calculating transformer flops mainly relies on the methodology outlined in Chinchilla~\citep{hoffmann2022training}, accounting for the flops associated with the feed-forward networks, the qkvo projections within self-attention, as well as both the attention computation and output projection steps. An important distinction in our analysis is the assumption that the input embedding layer utilizes a highly efficient lookup mechanism rather than being performed via dense matrix multiplication, and as such, we treat this step as requiring zero flops. Consistent with prior research, we also assume that the backward propagation requires double the flops of the forward propagation.

To compute flops \textit{per byte} for {\textsc BLT} models, we add up the flops for the local encoder transformer, the global latent transformer, and the local decoder transformer, together with the cross attention blocks in the encoder and the decoder:

\begin{align}
    \text{FL}_{\text{{\textsc BLT}}} &= \text{Transf. FL}(h_{\mathcal{G}}, l_{\mathcal{G}}, m=n_{ctx}/n_p, V=0) / n_p \\
    &+ \text{Transf. FL}(h_{\mathcal{E}}, l_{\mathcal{E}}, m=w_{\mathcal{E}}, V=0) \\
    &+ \text{Transf. FL}(h_{\mathcal{D}}, l_{\mathcal{D}}, m=w_{\mathcal{D}}, V=256) \\
    &+ \text{Cross Attn. FL}(h_{\mathcal{E}}, l_{\mathcal{E}}, m=n_p, r=n_p/k) \cdot k/n_p \\
    &+ \text{Cross Attn. FL}(h_{\mathcal{D}}, l_{\mathcal{D}}, m=k, r=k/n_p)
\end{align}

Here, $n_{ctx}$ denotes the sequence length in bytes, while $n_p$ represents the patch size. The parameter $r$ specifies the ratio of queries to key/values, and $k$ indicates the proportion of patch dimension to byte dimension—that is, it describes how many local model splits are concatenated to produce the global model representation ($k=2$ in \autoref{fig:crossattn}). The variable $V$ is the vocabulary size utilized for output projection, which is implemented exclusively within the local decoder. The specific context length, $m$, that the attention mechanism relies on, varies depending on whether the respective module is applied to byte sequences or patch sequences. For each module, we accordingly adapt the attention flop calculations. Precise formulations used for calculating FLOPs in Transformer and Cross-Attention modules are detailed in Appendix~\ref{section:flops-appendix}.

\subsection{Bits-Per-Byte Estimation} Perplexity only makes sense in the context of a fixed tokenizer as it is a measure of the uncertainty for each token. When comparing byte and token-level models, following previous work~\citep{xue2022byt5, yu2023megabyte, wang2024mambabyte}, we instead report Bits-Per-Byte (BPB), a tokenizer independent version of perplexity. Specifically:

\begin{align}
    \text{BPB}(x)&= \frac{\mathcal{L}_{CE}(\pmb{x})}{\ln(2)\cdot n_{\text{bytes}}}
\end{align}

In this formulation, the total cross-entropy loss computed for the data $\pmb{x}$ serves as an uncertainty metric, which is subsequently divided by both the overall count of bytes in $\pmb{x}$ and a normalization constant.

\subsection{Transformer Architecture Hyperparameters} Across all transformer layers in {\textsc BLT}, encompassing both local and global models, our configuration generally mirrors that of Llama~3~\citep{dubey2024llama}. Specifically, we employ the SwiGLU activation function~\citep{swiglu} within the feed-forward modules, integrate rotary positional embeddings (RoPE)~\citep{RoPe} with $\theta=500000$~\citep{xiong2024effective} exclusively in the self-attention mechanisms, and utilize RMSNorm \citep{RMSNorm} for normalization throughout the network. For all self-attention components utilizing static attention masks such as \textit{block causal} or \textit{fixed-window block causal}, Flash attention~\citep{dao2022flashattention} is adopted, with a standard window size of 512 for fixed-width mask configurations. In the context of cross-attention modules, where patch-dependent dynamic attention masks are required, Flex Attention\footnote{\url{https://pytorch.org/blog/flexattention}} is incorporated to enable efficient fused operations and substantially accelerate the training process.

\subsection{{\textsc BLT}-Specific Hyperparameters} To evaluate the performance of {\textsc BLT} models, we perform experiments focusing on two main aspects: scaling behavior and downstream task performance. For these assessments, we explore model sizes of 400M, 1B, 2B, 4B, and 8B parameters. The detailed architectural configurations for each scale are included in Appendix~Table~\ref{tab:arch}.

To initialize queries in the first cross-attention module of the local encoder, we employ max-pooling. Across all {\textsc BLT} variants, we apply $500,000$ hashes generated by a single hash function, utilizing n-gram lengths between 3 and 8. The learning rate is set to $4e-4$ for all scales. This selection is informed by a hyperparameter search over the interval $1e-3$ to $1e-4$ at the 400M and 1B parameter sizes, where results indicated that using the same learning rate as token-based models is optimal. Training on Llama-2 data follows the batch-size recommendations by~\cite{dubey2024llama}, implemented either by sample count or equivalent byte size.

For optimization, we choose AdamW~\citep{adamw} with hyperparameters $\beta_1=0.9$, $\beta_2=0.95$, and $\epsilon=10^{-8}$. The learning rate schedule incorporates a linear warm-up over 2000 steps, followed by cosine annealing to zero. We apply a weight decay coefficient of 0.1 and implement global gradient norm clipping at a maximum of 1.0.

\section{Scaling Trends}
\label{section:scaling}

We offer a comprehensive analysis of the scaling characteristics exhibited by byte-level models, with the intention of providing insights to guide the future scaling of {\textsc BLT} models. Our investigation seeks to overcome shortcomings observed in earlier studies of byte-level models by: (a) contrasting scaling patterns under compute-optimal training conditions, (b) training comparable 8B parameter models on large-scale datasets (reaching up to 1T tokens or 4T bytes) followed by downstream evaluations, and (c) examining scaling behavior within settings where inference costs are kept constant. Later in the paper, we further explore the unique strengths associated with modeling byte-sequences.

\subsection{Parameter Matched Compute Optimal Scaling Trends}
\label{section:parameter-matched-scaling}
We utilize the Llama 2 dataset to train a selection of \textit{compute-optimal} bpe and {\textsc BLT} models, covering four model sizes between 1B and 8B parameters. Performance is evaluated by plotting the required training flops against language modeling accuracy on a representative subset of the training mixture. The bpe models are trained following the optimal parameter-to-data ratio prescribed by Llama~3~\citep{dubey2024llama}, ensuring an efficient allocation of compute per the training dataset. This \textit{compute-optimal} regime is established to maximize performance under fixed computational resources~\citep{hoffmann2022training}, serving as a reliable point of comparison for evaluating other architectures. In parallel, for each bpe model, we train a {\textsc BLT} model of equivalent size and architectural configuration, utilizing a Latent Transformer and the identical training data.

As shown in~\autoref{fig:scaling-compute-opt} (right), {\textsc BLT} models consistently equal or surpass the performance of their bpe-based analogues, with this pattern persisting as both model capacity and computational requirements are increased. To our knowledge, {\textsc BLT} represents the first byte-level Transformer design to display comparable scaling properties to those of BPE-based architectures under compute-optimal conditions. These results support our hypothesis that the ideal proportion between parameters and training computation established for bpe remains relevant for {\textsc BLT}, or at the very least, does not significantly deviate from it.

To achieve bpe scaling trends, it is essential to utilize both advancements in architecture and dynamic patching strategies. In ~\autoref{fig:scaling-compute-opt} (left), we present a comparison between models employing space-patching techniques and Llama 3. For this evaluation, SpaceByte \citep{slagle2024spacebyte} is estimated using {\textsc BLT} space-patching that omits n-gram embeddings and cross-attention mechanisms. While SpaceByte yields improvements relative to Megabyte, its performance still falls significantly behind Llama 3. On the right side of \autoref{fig:scaling-compute-opt}, the combined effects of architectural refinements and dynamic patching are depicted. At scale, {\textsc BLT} models demonstrate performance on par with the leading tokenizer-based models, such as Llama 3.

Additionally, our results indicate that, for tokenizer-based models, the selection of tokenizer notably impacts model performance; specifically, models utilizing the Llama-3 tokenizer achieve higher accuracy than those employing the Llama-2 tokenizer when trained on identical datasets.

In summary, the {\textsc BLT} architecture demonstrates a performance pattern intermediate between Llama 2 and Llama 3 when adopting substantially larger patch sizes. While the bpe tokenizers utilized in Llama 2 and 3 exhibit mean token sizes of 3.7 and 4.4 bytes respectively, {\textsc BLT} achieves comparable scaling characteristics at average patch sizes of 6 or even 8 bytes. Since inference flop decreases as the average patch size increases, an 8-byte patch size yields almost 50\% reduction in inference flops. Furthermore, scaling both model and dataset size appears to favor models employing larger patch sizes. For instance, {\textsc BLT} configured with an 8-byte patch size initially underperforms relative to the bpe version of Llama 2 at the 1B parameter level, but surpasses its bpe counterpart when scaled up to 7B parameters. This pattern indicates that such larger patch sizes could be increasingly advantageous at greater model scales, and it raises the possibility that even more extensive patch sizes may be viable as both model capacity and computational resources are expanded.

\subsection{Beyond Compute Optimal Task Evaluations}
\label{section:evals}
To further investigate scaling behavior, we train an 8B {\textsc BLT} model at a scale exceeding the compute optimal ratio using the {\textsc BLT}-1T dataset, which is larger and of higher quality. We then evaluate the resulting model on various established classification and generation benchmarks. For this evaluation, we focus on several standard tasks targeting common sense reasoning, general world knowledge, and code generation:
\paragraph{Classification tasks} considered are ARC-Easy (0-shot)~\citep{Eval_ARC}, Arc-Challenge (0-shot)~\citep{Eval_ARC}, HellaSwag (0-shot)~\citep{Eval_hellaswag}, PIQA (0-shot)~\citep{Eval_piqa}, and MMLU (5-shot)~\citep{hendrycks2020measuring}. We utilize prompt-scoring, where probabilities are assigned to possible answer choices, and performance is summarized as the mean accuracy across tasks.
\paragraph{Coding related generation tasks:}  For assessing Python code synthesis, we present pass@1 outcomes on MBPP (3-shot)~\citep{Eval_MBPP} and HumanEval (0-shot)~\citep{Eval_HumanEval}, thereby gauging the LLMs’ code generation capabilities.

In \autoref{tab:evals}, we present a comparative analysis of three models trained using the {\textsc BLT}-1T dataset: a model based on the Llama 3 bpe tokenizer,\footnote{Our selection of the Llama 3 tokenizer, which has a 128k vocabulary, is motivated by its superior performance relative to the Llama 2’s 32k vocabulary.} and two variants of the {\textsc BLT} approach. One variant employs a space-patching methodology ({\textsc BLT}-Space), while the other implements an entropy-driven patching strategy ({\textsc BLT}-Entropy). In the entropy-based approach, we apply an approximate monotonicity constraint and refresh the model’s context at each newline, as detailed in~\autoref{section:entropy-context}. All models are trained under comparable flop budgets to ensure fair evaluation. Notably, for {\textsc BLT}-Entropy, we also reduce the entropy threshold during inference from 0.6 to 0.1, a modification that enhances task performance though it necessitates a greater number of inference steps.

The {\textsc BLT}-Entropy model achieves superior performance compared to the Llama 3 model on 4 out of 7 evaluated tasks, despite both being trained on an equivalent number of bytes. This enhancement in results can be attributed to two main factors: (1) more efficient utilization of training compute enabled by dynamic patching, and (2) explicit modeling of information at the byte level rather than at the token level.

Conversely, {\textsc BLT}-Space demonstrates inferior performance relative to the Llama 3 tokenizer in all tasks except one; however, its utilization of a larger mean patch size of 6 bytes leads to a notable decrease in inference flops. In contrast, models utilizing bpe and entropy-patching maintain a similar average patch size, close to 4.5 bytes, across the training data mix. Under an identical training budget, the model with the larger patch size is able to process 30\% more data than the other two, potentially resulting in {\textsc BLT} deviating further from compute-optimal efficiency.

\subsection{Patches Scale Better Than Tokens} {\textsc BLT} models provide the flexibility to enlarge both the model and the patch sizes simultaneously without exceeding a fixed training and inference computational budget, or increasing the quantity of training data. The ability to scale patch size at will distinguishes patch-based approaches from traditional fixed-vocabulary token-based models, which face inherent efficiency constraints, as described in Section~\ref{section:bpe-patch}. Using larger patch sizes reduces the total computational load, allowing more resources to be allocated toward scaling up the global latent transformer, since it needs to be executed less frequently.

We carry out a fixed inference scaling experiment to evaluate whether increasing model size—paired with fewer processing steps over larger patches—yields improved performance relative to smaller models executing more steps. Our analysis begins with models of 400 million and 3.6 billion parameters, utilizing the Llama 2 tokenizer, and identifies flop-equivalent counterparts using the Llama 3 tokenizer as well as {\textsc BLT}-Entropy models. These {\textsc BLT}-Entropy models have average patch sizes of 6 and 8 bytes when trained on the designated data mixture (for further specifics, refer to~\autoref{tab:fixed-inf-params}). For models with patch size 8, we incorporate 3 encoder layers in lieu of just one. Each configuration is trained across different training flop budgets.

As illustrated in \autoref{fig:fixed_inference_scaling}, {\textsc BLT} models demonstrate superior scaling properties compared to tokenization-based approaches across both inference flop categories. While BPE models initially outperform at lower training compute allocations, {\textsc BLT} soon overtakes them shortly after exceeding the compute-optimal region. In real-world scenarios, it is often advantageous to invest additional resources into pretraining, since this leads to higher-performing models even when constrained by a fixed inference budget. This pattern can be observed in the case of 8B parameter models such as Llama 3.1, which was pretrained on a dataset size approximately two orders of magnitude larger than that suggested by compute-optimality for its scale.

The point at which {\textsc BLT} surpasses token-based models has moved marginally closer to the compute-optimal threshold in the case of the larger flop class models, shifting from 3x to 2.5x of the compute-optimal budget. Likewise, the model with patch size 8 exhibits a more rapid scaling behavior in the higher flop regime, overtaking the other configurations at an earlier stage. As outlined in Section~\ref{section:parameter-matched-scaling}, increasing the patch size tends to bring model performance more in line with BPE models as model scale grows. We partly ascribe this phenomenon to the fact that the relative contribution of the byte-level Encoder and Decoder modules to total flops diminishes, since their computational requirements scale less sharply than those of the Latent Transformer. Notably, when the total parameter count expands 20-fold from 400M to 8B, the number of local parameters specific to {\textsc BLT} increases by only about twofold. This distinction is critical because larger patch sizes exclusively influence the computational load associated with the patch Latent Transformer, leaving the byte-level modules largely unaffected. As a result, the {\textsc BLT}-Entropy ps=8 model's parameter count rose from 1.6x to 1.7x that of the Llama 2 model upon scaling to a larger model size.

In conclusion, our investigation into patch-length scaling reveals that the {\textsc BLT} patch-based framework exhibits superior scaling behavior when both patch size and model size are increased together. These positive scaling patterns not only continue but can even strengthen as model sizes grow larger.

\section{Byte Modeling Improves Robustness} Additionally, we assess the robustness of {\textsc BLT} relative to models that operate solely on tokens without incorporating byte-level cues. We further introduce a method for augmenting pretrained token-based models with byte-level capabilities.
\label{sec:robustness}

\subsection{Character-Level Tasks}

One of the initial incentives for developing byte-level models was their inherent resilience to noise at the byte level in the input, along with their ability to recognize sub-token elements—an area where models based on tokenization often face limitations. To systematically assess these properties, we conduct supplementary experiments using benchmark datasets designed to test both robustness against noisy inputs and sensitivity to character-level details. These evaluations span English and multiple other languages, covering not only alphabetic characters but also numerals and phonetic symbols. The findings from these evaluations are summarized in Table~\ref{tab:char_tasks}.

\paragraph{Noisy Data} To assess the resilience of {\textsc BLT} compared to tokenizer-based architectures, we generate corrupted variants of the standard classification datasets outlined in Section~\ref{section:evals}. Five different character-level corruption techniques are employed to inject diversity into the textual input: (a) \textit{AntSpeak}: The entire text is transformed into uppercase letters separated by spaces. (b) \textit{Drop}: 10\% of all characters are randomly deleted from the input. (c) \textit{RandomCase}: Randomly alternates 50\% of characters to uppercase and 50\% to lowercase. (d) \textit{Repeat}: 20\% of characters are selected and each is duplicated up to four times. (e) \textit{UpperCase}: All characters in the input are rendered in uppercase. During testing, each corruption method is independently applied to the prompt, the completion, or to both segments, treating each setting as a distinct task; we then compute and report the mean performance. The outcomes, summarized in Table~\ref{tab:char_tasks} on noised HellaSwag~\citep{Eval_hellaswag}, demonstrate that {\textsc BLT} consistently exhibits greater robustness relative to tokenizer-based systems, exceeding models trained on equivalent data by an average of 8 points, and surpassing the Llama 3.1 model which leverages a substantially larger training corpus.

\paragraph{Phonology - Grapheme-to-Phoneme (G2P)} We evaluate how effectively {\textsc BLT} translates sequences of graphemes—written symbols comprising a word—into their corresponding phonemic representations, which capture the word’s spoken form. Table \ref{tab:char_tasks} displays the G2P task outcomes within a 5-shot learning framework, employing the Phonology Bench~\citep{suvarna2024phonologybench}. Our analysis indicates that {\textsc BLT} achieves superior performance in comparison to the Llama 3 1T tokenizer-based model baseline for this task.

\paragraph{CUTE}
To evaluate character-level comprehension, we test {\textsc BLT} on the CUTE benchmark~\citep{edman2024cute}, which contains a range of tasks grouped into three main types: compositional understanding, orthographic similarity, and sequence manipulation capabilities. This benchmark presents a formidable challenge for most models relying on tokenization, as these typically recognize how their tokens are spelled, yet encounter difficulties leveraging this knowledge to transform or alter text accurately. As shown in Table~\ref{tab:char_tasks}, {\textsc BLT}-Entropy surpasses both BPE Llama 3 models by over 25 points on this benchmark. Notably, the model exhibits outstanding results on character manipulation tasks, scoring 99.9\% accuracy on both spelling-focused tests. These substantial gains—despite {\textsc BLT} being trained with just one-sixteenth the data used by Llama 3.1—suggest that mastering character-level information poses a greater challenge for BPE-based models. Figure \ref{fig:cute_examples} displays several instances in which the Llama 3 tokenizer model falters whereas our {\textsc BLT} model succeeds. The only categories where the BPE approach shows an advantage are word deletion and insertion. These kinds of word-level modifications are potentially less accessible to a byte-level model, yet the performance differential is not pronounced, and constructing word representations from characters may, in fact, prove easier than the inverse. Our experiments employ a consistent evaluation framework and the original prompts provided by Huggingface across all tasks. Additional prompt engineering may yield further improvements for BPE models.

\paragraph{Low Resource Machine Translation} We assess {\textsc BLT} on translation tasks involving twenty-one lower-resourced languages from six major language families, spanning diverse scripts, using the FLORES-101 benchmark~\citep{goyal2022flores}. SentencePiece BLEU results are summarized in Table~\ref{tab:flores}. Our findings indicate that {\textsc BLT} achieves superior performance compared to a baseline model employing the Llama 3 tokenizer, securing an overall increase of 2 BLEU points for English-target translation and an improvement of 0.5 BLEU points for English-source translation. On well-studied language pairs, {\textsc BLT} achieves results that are on par with or marginally surpass those of Llama 3. Notably, for many pairs in low-resource language families, {\textsc BLT} consistently exceeds the performance of Llama 3, highlighting the value of byte modeling in handling long-tail byte sequence generalization.

\subsection{Training {\textsc BLT} from Llama 3}
\label{sec:distillation}
We investigate a training approach where {\textsc BLT} benefits from pre-trained tokenizer-based models to achieve improved and accelerated convergence. Specifically, we initialize the global transformer parameters of {\textsc BLT} using weights from a pre-trained Llama 3.1 model. Following this, during training, we apply a learning rate for the global transformer that is one-tenth of the rate used for the local encoder and local decoder, aligning the number of optimization steps to the standard used for Llama 3. A comparative analysis against a baseline {\textsc BLT} is provided in Table~\ref{tab:distill}. The results demonstrate that initializing {\textsc BLT} from Llama 3.1 leads to substantial performance improvements over both the original Llama 3 and standard {\textsc BLT}, under equal computational budgets. Furthermore, relative to our {\textsc BLT}-Entropy configuration from Table~\ref{tab:evals}, which was trained with a much larger dataset (1T tokens), the {\textsc BLT} initialized from Llama 3.1 continues to deliver enhanced results on the MMLU benchmark. This indicates that such initialization can notably reduce training flops without compromising on performance.

This approach can alternatively be interpreted as adapting models that rely on tokenizers into ones that operate without them, in effect transforming a pre-trained LLaMA 3.1 model into a {\textsc BLT} model. For a thorough evaluation, Table \ref{tab:distill} presents both the original LLaMA 3.1 (trained on 15T tokens) and the {\textsc BLT} version initialized from LLaMA 3, and compares their results. While our model shows minor reductions in accuracy on MMLU and HumanEval, its performance drops more substantially on a range of other benchmarks. These findings indicate the necessity of additional research to fully harness the strengths of the pre-trained model—especially with respect to optimizing the composition of training data and fine-tuning hyperparameters.

\section{Ablations and Discussion}
\label{section:ablations}
Here, we present a series of ablation studies that validate the architectural decisions made for {\textsc BLT}, as well as the selection of patching strategy and hyper-parameter configuration used for the {\textsc BLT} 8B model, which was trained with the {\textsc BLT}-1T dataset.

\paragraph{Entropy Model Hyper-parameters} To investigate how adjustments in entropy model capacity and context window size influence scaling behavior, we develop byte-level entropy transformer models varying in size from 1m to 100m parameters, and utilize context window lengths ranging from 64 up to 512. Bpb as a function of training flops is visualized in scaling law curves, which are generated for our $400m$ and $1b$ {\textsc BLT} models—both trained on the Llama-2 dataset—and are illustrated in \autoref{fig:entropy_ablation}. The results indicate that increasing either model size or context window enhances scaling, although the benefits plateau after about 50m parameters.

\paragraph{Types of Patching}

We conduct ablation studies on the four patching methods outlined in Section~\ref{section:patching}, namely: 1) Strided Patching with strides of 4 and 6, 2) Whitespace-based Patching, 3) BPE Tokenizer patching employing the Llama 3 tokenizer, and 4) Entropy-driven patching utilizing a compact byte-level language model.

Although dynamic patching decreases the actual sequence length, we ensure comparable context lengths across different patching strategies by holding the sequence length constant. This design means that, on average, every model processes an equivalent number of bytes per sequence throughout training and inference, thereby mitigating confounding influences linked to modeling longer contexts. The outcomes of these ablation studies are summarized in Figure~\ref{fig:scaling-compute-opt}. Notably, all alternative patching strategies surpass static patching, and space patching performs nearly as well as dynamic entropy-based patching.

\autoref{tab:evals_ablation} reports benchmark results for {\textsc BLT} models, comparing approaches that utilize standard tokenizers, space patching, and entropy-based patching. All models are trained on the Llama 2 dataset for the optimal number of steps as identified in \citet{dubey2024llama}. While space patching offers a more straightforward alternative, as it eliminates the need to evaluate entropy models dynamically during training, our results confirm that the improvements seen with entropy-based patching on scaling properties (see Section~\ref{section:scaling}) consistently extend to performance on downstream benchmark tasks.\footnote{Space patching evaluations are drawn from earlier experiments conducted without cross-attention, but comparable patterns emerge even in the presence of cross-attention.}

\paragraph{Cross-Attention} 
\autoref{tab:crossattn_ablations_1b} presents an ablation study on the effects of introducing cross-attention modules at different locations within both the encoder and decoder components of {\textsc BLT}. Regarding encoder cross-attention, we evaluate three strategies for initializing the queries: (1) employing a uniform learned embedding shared across all global states; (2) utilizing a hash embedding derived from the patch bytes; and (3) applying a pooling operation over the encoder's hidden state representations of the patch bytes at the corresponding encoder layer.

Our results indicate that incorporating cross-attention within the \textit{decoder} yields the highest effectiveness. Within the encoder, the benefits of cross-attention are minimal and are observed solely when queries are initialized using pooling. Furthermore, cross-attention proves especially advantageous for Common-Crawl data, with the effect being more pronounced for larger patch sizes.

\paragraph{n-gram Hash Embeddings} 

We conduct an ablation study by testing n-gram hash embedding vocabulary sizes of 0, 100K, 200K, and 400K, with the corresponding results summarized in Table~\ref{tab:ngram_ablations_1b}. Our findings indicate that hash embeddings consistently enhance performance across all domains, with the most pronounced effects observed on Wikipedia and Github—showing a 0.04 bpb improvement compared to a 0.01 bpb gain after 15k steps at the 8B scale. Furthermore, reducing the number of hashes from 500K to 300K at the 8B scale only led to an approximate decrease of 0.001 bpb after 15k steps, suggesting that hash embeddings are crucial for narrowing the performance gap between {\textsc BLT} and tokenizer-based models. Nevertheless, we note that benefits plateau once 300K hashes are used, demonstrating diminishing returns. We also observe that the advantages provided by hash embeddings and cross-attention tend to supplement each other, as they offer benefits across distinct datasets.

\paragraph{Local Model Hyperparamaters} Table \ref{tab:local_layers} presents an ablation study exploring different configurations for the number of layers in both the local encoder and decoder. Our results show that, when using hash n-gram embeddings, {\textsc BLT} achieves strong performance with a highly simplified encoder—specifically, one comprising only a single layer—coupled with a more complex, deeper decoder.

\section{Related Work}
\label{section:related_work}

\textbf{Character-Level RNNs:} Since the advent of neural networks, Character Language Modeling has garnered significant attention~\citep{sutskever2011generating,mikolov2012subword,graves2013generating}, particularly due to its inherent capability to seamlessly handle out-of-vocabulary words without relying on back-off approaches. In addition, \cite{kim2016character} developed a model that exclusively processes input as character sequences via convolutional and highway networks, subsequently connecting these representations to LSTM-based RNNs. This architecture achieved parity with leading RNN-based language models of that era in English and exceeded their performance on languages characterized by rich morphology—highlighting a valuable benefit of character-level LLMs. Expanding on this concept, \cite{kenter2018byte} utilized byte-level LSTM architectures for machine comprehension tasks, demonstrating improvements over word-based counterparts on morphologically complex languages such as Turkish and Russian. Similarly, \cite{zhang2015character} explored character-driven convolutional approaches in classification settings, and for certain applications, these surpassed traditional word-level techniques. In a different approach, \citet{chung2019hierarchical} proposed hierarchical LSTM architectures, equipping each layer with boundary detectors to uncover latent hierarchical structures within the text, thus bolstering the effectiveness of character-level language models. Departing from attention-based frameworks, ByteNet by \cite{kalchbrenner2016neural} introduced character-level processing using CNN layers for the purpose of machine translation.

\textbf{Character-Level Transformers:} The introduction of attention-based transformer architectures~\citep{vaswani2017attention}, in combination with subword tokenization techniques~\citep{sennrich-etal-2016-neural}, resulted in notable gains for neural language modeling and performance on various benchmarks. Nevertheless, employing word or sub-word tokens inherently imposes an inductive bias on the abstraction layer at which models process language. In an effort to merge the achievements of transformers with encouraging early findings from character-based language models, \citet{al2019character} implemented deep transformer networks. By incorporating auxiliary loss functions, they were able to train transformer models that surpassed prior LSTM approaches for character-level language modeling. Despite this, their results still lagged behind those achieved by word-level large language models. Similarly, GPT-2~\citep{radfordgpt2} noted that when trained on extensive datasets, such as the 1 Billion Word Benchmark, byte-level language models did not match the effectiveness of their word-level counterparts.

Although \cite{Choe2019BridgingTG} showed that transformer-based LLMs operating at the byte level can exceed the performance of comparably sized subword-level LLMs, these models come with significantly higher computational costs and prolonged training durations. Likewise, \cite{el2020characterbert} introduced CharFormer, a BERT-based architecture that leverages convolutions on character embeddings to form word representations, demonstrating gains in the medical domain but incurring substantial computational overhead. In a related effort, \cite{clark2022canine} presented CANINE, an encoder-only architecture with 150M parameters that processes sequences at the character level. CANINE combines a deep transformer backbone, comparable to our global module, with localized transformers and strided convolutions for input character downsampling. This design allows CANINE to surpass the corresponding token-level encoder-only baseline (mBERT) on multilingual tasks. The ByT5 model~\citep{xue2022byt5} evaluated byte-level encoder-decoder systems that eliminate patching operations entirely. Despite showing heightened resilience to input noise and competitive results compared to token-based models (with only a quarter of the data), ByT5’s approach—eschewing patching—resulted in high computational demands, as expensive attention was required over every single byte. Modeling raw bytes instead of subword tokens inherently leads to longer input sequences, which compounds scaling challenges for byte-level models. More recently, leveraging the Mamba Architecture~\citep{gu2023mamba}, which supports constant-size memory states over extensive contexts, \cite{wang2024mambabyte} trained a byte-level version of Mamba—again, without patching—demonstrating superior performance relative to byte-level transformer counterparts at the 350M parameter scale, as measured by bits-per-byte across diverse datasets in flop-equivalent scenarios.

\textbf{Patching-based approaches:} Leveraging patching strategies can substantially reduce the computational overhead (measured in flops) typically associated with byte-level LLMs, all while maintaining or only modestly impacting model effectiveness. Several studies have shown promising outcomes, particularly with relatively smaller models and modest training dataset sizes. \cite{nawrot-etal-2022-hierarchical} implemented static patching that involves downsampling and upsampling within the hourglass transformer architecture, achieving superior performance compared to alternative byte-level benchmarks at the 150M parameter scale. Building on this, \cite{nawrot-etal-2023-efficient} introduced dynamic patching techniques, incorporating innovations such as a boundary-predictor (both trained end-to-end and guided by certain tokenizers) and an entropy-informed patching mechanism akin to {\textsc BLT}. Their results indicate that these methods surpass the standard transformers of the era for language modeling at the 40M parameter level with training on approximately 400M tokens. Additionally, \cite{lester2024training} explored language model training over sequences that had been compressed with arithmetic coding to exceed the compression efficiency of BPE. Utilizing an equal-info windows method, their models eclipsed byte-level baselines on language modeling benchmarks, though they remained less effective than models trained on subword-level data.

Our study builds upon and is closely aligned with the MegaByte approach~\citep{yu2023megabyte}, a decoder-only causal large language model that employs a predetermined, unchanging method for partitioning bytes into patches and concatenating representations. It further applies a local decoding model to reconstruct bytes from these patches. MegaByte demonstrates competitive performance with tokenizer-based architectures at a parameter count of 1B on datasets containing approximately 400B bytes. Across our investigations, we systematically compare to MegaByte, observing that the method of fixed static patching falls short of matching the latest compute-efficient tokenizer-based models, particularly in settings where the computation budget (FLOPs) is controlled. Our results highlight how {\textsc BLT} effectively narrows this performance disparity. Parallel findings by \cite{slagle2024spacebyte} note similar limitations in MegaByte, proposing enhancements that include adapting the patching strategy to occur at whitespaces or other space-representative bytes and incorporating an additional local encoder module. Their results demonstrate gains over tokenized transformer models in constrained computation environments for specific domains, such as Github and arXiv, at the 1B parameter scale. We also present experimental results on their model and conclude that advancing byte-level models to truly rival cutting-edge token-based systems like Llama 3 necessitates further architectural advancements.

\section{Limitations and Future Work}

For this study, in making architectural decisions, we trained models for the ideal number of steps as established for Llama~3~\citep{dubey2024llama}. Nonetheless, since these scaling laws were derived for BPE-level transformers, applying them to {\textsc BLT} could yield less-than-optimal ratios between data and parameter count. Determining scaling laws specific to {\textsc BLT}, which may reveal even more advantageous scaling characteristics for our design, is left for future research. Furthermore, since a majority of our experiments were limited to models up to 1B parameters, it remains possible that the optimal configurations will shift as we extend to 8B parameters and larger, potentially resulting in enhanced performance at increased scale.

Current transformer libraries and repositories are primarily optimized for the efficiency of tokenizer-based transformer models. Although we report experiments matched in theoretical floating point operations and utilize some optimized methods (including FlexAttention) to accommodate layers that differ from standard transformer structures, our practical implementations may still fall short of the runtime performance seen in tokenizer-based models and could see improvements through additional optimization efforts.

Whereas {\textsc BLT} relies on an independently trained entropy model for patching, integrating the patching model into an end-to-end learning framework presents a promising avenue for future investigation. As described in Section \ref{sec:distillation}, we provide preliminary results that suggest the viability of ``byte-ifying'' tokenizer-based architectures such as Llama 3—particularly those trained on datasets exceeding 10T tokens—by initializing and freezing the global transformer using pretrained weights. Continued exploration of this approach could yield strategies that not only maintain the advantages afforded by byte-level modeling, but also achieve superior performance compared to these tokenizer-based models, all without the need for complete retraining.

\section{Conclusion}
\label{section:conclusion}

In this work, we introduce the Byte Latent Transformer (\textbf{BLT}), an innovative model architecture that challenges the standard reliance on fixed-vocabulary tokenization found in contemporary large language models. {\textsc BLT} departs from traditional approaches by employing a flexible, end-to-end learnable grouping mechanism that aggregates bytes into adaptive patches, thereby tailoring computational effort to match the underlying complexity of the input. This methodological shift yields notable gains in both computational efficiency and model resilience. Through comprehensive scaling experiments, we show that {\textsc BLT} is competitive with tokenization-centric architectures, such as Llama 3, at scales reaching 8B parameters and 4T bytes; moreover, it can achieve up to a 50\% reduction in inference flops with only marginal sacrifices in key evaluation scores. Additionally, the {\textsc BLT} architecture introduces an alternative scaling strategy: enabling concurrent expansion of both the model and patch dimensions while maintaining a constant inference cost—a property especially beneficial in real-world computational constraints. Crucially, by processing data directly at the byte level, {\textsc BLT} demonstrates enhanced robustness when confronted with noisy or rare inputs and a richer capacity for modeling fine-grained sub-word information. Collectively, these findings establish {\textsc BLT} as a compelling contender to token-based methods, offering a robust and highly scalable path towards more efficient and versatile language modeling.

\end{document}